% =============================================================================
% Kapitel 3: Die Software EvAX
% =============================================================================
\chapter{Die Software EvAX}
\label{ch:evax}

\section{Überblick und Einordnung}

\evax{} (Evolutionary Algorithm for XAS) ist ein Programmpaket zur
Analyse von EXAFS-Daten, das an der Universität Lettland entwickelt
wurde \cite{Timoshenko2009, Timoshenko2012a, Timoshenko2012b}.
Im Unterschied zu konventionellen Least-Squares-Verfahren
-- wie sie etwa in \textsc{Artemis}/\textsc{Demeter}
\cite{Ravel2005} implementiert sind -- arbeitet \evax{} nicht mit
einer parametrischen Beschreibung der radialen Verteilungsfunktionen,
sondern optimiert die atomaren Positionen in einem dreidimensionalen
Strukturmodell direkt.

Die zentrale Idee besteht darin, eine Superzelle mit einigen hundert
Atomen aufzubauen, für jeden Absorbertyp das zugehörige EXAFS-Signal
ab initio zu berechnen und die Übereinstimmung mit dem Experiment
durch iterative Verrückung einzelner Atome zu verbessern.
Die Kombination eines evolutionären Algorithmus mit Elementen der
Reverse-Monte-Carlo-Methode (RMC) stellt sicher, dass der
Konfigurationsraum effizient abgesucht wird und lokale Minima
überwunden werden können.

\section{Grundprinzip: RMC und evolutionärer Algorithmus}

\subsection{Reverse Monte Carlo (RMC)}

Die RMC-Methode \cite{McGreevy1988} geht von einem atomistischen
Startmodell aus und modifiziert es durch zufällige Verschiebungen
einzelner Atome.
Nach jeder Verrückung wird das resultierende Signal berechnet und
mit dem Experiment verglichen.
Verbessert sich die Übereinstimmung (gemessen durch ein
Residuenfunktional), wird die Verschiebung akzeptiert; verschlechtert
sie sich, wird die Verschiebung mit einer bestimmten
Wahrscheinlichkeit dennoch zugelassen -- analog zum
Metropolis-Kriterium in der statistischen Physik.
Dieser stochastische Akzeptanzmechanismus verhindert, dass der
Algorithmus in einem lokalen Minimum verharrt.

\subsection{Evolutionärer Algorithmus}

\evax{} erweitert das RMC-Prinzip um einen evolutionären Algorithmus:
Es werden nicht eine, sondern $N_{\text{states}}$ (typischerweise 16)
parallele Strukturkopien (\emph{States}) gleichzeitig optimiert.
In regelmäßigen Abständen werden die States anhand ihrer Fitgüte
bewertet; schlechte Strukturen werden durch Kopien guter Strukturen
ersetzt, wobei leichte Mutationen eingeführt werden.
Dieses Selektions-Mutations-Schema entspricht dem Prinzip der natürlichen
Auslese und erhöht die Wahrscheinlichkeit, das globale Minimum zu finden
\cite{Timoshenko2014}.

\section{FEFF-Integration und Spektrumberechnung}

Für die Berechnung des theoretischen EXAFS-Signals verwendet \evax{}
das Programm \feff{} (Version~8.5) \cite{Rehr2010} als externen
Rechenkern.
\feff{} bestimmt die Streuphasen und -amplituden ab initio aus der
Muffin-Tin-Approximation des Kristallpotentials und erzeugt die
Beiträge aller relevanten Einfach- und Mehrfachstreupfade.

Da eine vollständige \feff{}-Rechnung für alle 256~Atome bei jeder
Iteration viel zu aufwendig wäre, nutzt \evax{} eine
Clusterexpansion: Die lokale Umgebung jedes Absorberatoms wird mit
einer Bibliothek bereits berechneter Cluster verglichen, und nur bei
hinreichend starker Veränderung wird eine neue \feff{}-Rechnung
ausgelöst.
Der Parameter \texttt{Update\_basis\_every} steuert, nach wie vielen
Iterationen die Streubasis aktualisiert wird.

\section{Multi-Edge-Auswertung}

Ein besonderer Vorteil von \evax{} liegt in der simultanen Auswertung
mehrerer Absorptionskanten.
Bei einer ternären Legierung wie CoNiCu werden die K-Kanten aller drei
Elemente gleichzeitig gefittet.
Da alle Kanten vom selben Strukturmodell erzeugt werden, sind die
Fitparameter implizit verknüpft: Eine Verrückung eines Co-Atoms
beeinflusst sowohl das Co-K- als auch das Ni-K- und Cu-K-Spektrum.
Dies erhöht die effektive Informationsdichte erheblich und reduziert
die Mehrdeutigkeit des Fits \cite{Timoshenko2012a}.

Die relative Gewichtung der einzelnen Kanten im Residuenfunktional wird
durch \evax{} automatisch angepasst (\texttt{adjust\_spectra\_weights}),
sodass keine Kante den Fit dominiert.

\section{Wichtige Eingabeparameter}
\label{sec:evax-params}

Tabelle~\ref{tab:evax-params} fasst die zentralen Parameter zusammen,
die in der Eingabedatei (\texttt{parameters.dat}) spezifiziert werden.

\begin{table}[htbp]
  \centering
  \caption{Zentrale Eingabeparameter von \evax{} und ihre Bedeutung.}
  \label{tab:evax-params}
  \small
  \begin{tabularx}{\textwidth}{lX}
    \toprule
    \textbf{Parameter} & \textbf{Bedeutung} \\
    \midrule
    \texttt{File\_with\_structure}
      & Pfad zur Startstruktur im POSCAR-Format (Gittervektoren und
        Atompositionen) \\
    \texttt{Supercell}
      & Multiplikation der Einheitszelle, z.\,B.\ \texttt{1x1x1} für
        eine bereits aufgebaute Superzelle \\
    \texttt{S02}
      & Amplitudenreduktionsfaktor $S_0^2$ für jede Kante.
        Wird mit der \texttt{-autoES}-Option automatisch bestimmt \\
    \texttt{dE0}
      & Energieverschiebung $\Delta E_0$ (\si{\electronvolt}) zwischen
        theoretischer und experimenteller Kantenlage \\
    \texttt{k\_min}, \texttt{k\_max}
      & Grenzen des genutzten $k$-Bereichs (\si{\per\Angstrom}) \\
    \texttt{k\_power}
      & Exponent $n$ der $k$-Gewichtung ($k^n \cdot \chi(k)$) \\
    \texttt{Space}
      & Vergleichsraum: \texttt{k} ($k$-Raum), \texttt{R} ($R$-Raum)
        oder \texttt{w} (Wavelet) \\
    \texttt{N\_legs}
      & Maximale Zahl der Streuarme pro Pfad (2\,=\,nur
        Einfachstreuung, 4\,=\,inkl.\ MS) \\
    \texttt{R\_max\_for\_FEFF}
      & Maximaler Radius für die \feff{}-Cluster (\si{\Angstrom}) \\
    \texttt{Maximal\_step\_length}
      & Maximale Verschiebung pro Atom und Iteration (\si{\Angstrom}) \\
    \texttt{Maximal\_displacement}
      & Maximale kumulative Verschiebung eines Atoms gegenüber der
        Startposition (\si{\Angstrom}) \\
    \texttt{Number\_of\_states}
      & Anzahl paralleler Strukturkopien im EA \\
    \texttt{Stop\_after}
      & Zahl der Iterationen in der Screening-Phase \\
    \texttt{Froze\_in}
      & Abbruchkriterium: Iterationen ohne Verbesserung bis zur
        Konvergenz \\
    \bottomrule
  \end{tabularx}
\end{table}

\section{Ausgaben und Ergebnisgrößen}

Nach Abschluss der Optimierung liefert \evax{} folgende Ergebnisse:

\begin{description}[style=nextline]
  \item[Finale Struktur (\texttt{final.xyz})]
    Kartesische Koordinaten aller Atome in der optimierten Superzelle,
    aus denen Paarverteilungsfunktionen, Koordinationszahlen und SRO-Parameter
    berechnet werden können.
  \item[EXAFS-Spektren (\texttt{EXAFS.dat})]
    Berechnetes und experimentelles $\chi(k)$ für jede Kante, zur
    visuellen und quantitativen Beurteilung der Fitgüte.
  \item[Fouriertransformierte (\texttt{expFT.dat}, \texttt{expBFT.dat})]
    FT des experimentellen und berechneten Signals.
  \item[Radiale Verteilungsfunktionen (RDF)]
    Elementaufgelöste $g_{ij}(R)$ für alle Paarkombinationen.
  \item[Konvergenzverlauf (\texttt{output.dat})]
    Residuen als Funktion der Iterationszahl für jede Kante und als
    gewichtete Summe.
\end{description}

\section{Beispielhafter Workflow in EvAX}
\label{sec:workflow}

Im Folgenden wird der typische Arbeitsablauf bei einer \evax{}-Analyse
Schritt für Schritt beschrieben.
Abbildung~\ref{fig:workflow} zeigt eine schematische Darstellung.

\subsection{Datenvorverarbeitung}

Die experimentellen Rohdaten liegen als Absorptionsspektren $\mu(E)$
vor und müssen zunächst in $\chi(k)$-Daten umgewandelt werden.
Dies geschieht in der Regel mit der Software \textsc{Athena}
\cite{Ravel2005} oder mit der Python-Bibliothek \textsc{larch}
\cite{Newville2013} und umfasst:
\begin{enumerate}[nosep]
  \item Bestimmung der Kantenenergie $E_0$,
  \item Abzug des Vor-Kanten-Hintergrunds (Pre-Edge-Subtraktion),
  \item Normierung auf den Kantensprung $\Delta\mu_0$,
  \item Entfernung des atomaren Hintergrunds mittels Spline-Fit
        (\texttt{autobk}-Algorithmus),
  \item Export der $\chi(k)$-Daten.
\end{enumerate}

Bei Multi-Edge-Messungen, bei denen mehrere Kanten in einem einzigen
Energiescan abgedeckt werden, muss zusätzlich darauf geachtet werden,
die Energiebereiche sauber zu separieren, damit die Normierung einer
Kante nicht durch die nächstfolgende verfälscht wird.

\subsection{Aufbau des Startmodells}

Das Strukturmodell wird als Superzelle im POSCAR-Format vorgegeben.
Typisch ist eine $4\times4\times4$-Vervielfachung der konventionellen
Einheitszelle, was für eine FCC-Struktur $4^3 \times 4 = 256$~Atome
ergibt.
Die chemischen Elemente werden zufällig auf die Gitterplätze verteilt,
wobei die Stöchiometrie der Probe eingehalten wird.

Die Wahl der Kristallstruktur (FCC, BCC, HCP) definiert die Topologie
des Startmodells.
\evax{} kann zwar Atome verschieben, aber nicht die grundlegende
Koordination des Gitters ändern.
Daher ist es sinnvoll, mehrere Ausgangsstrukturen zu testen und
anhand der Fitresiduen zu entscheiden, welche Struktur die Daten
am besten beschreibt.

\subsection{Konfiguration und Start}

Alle Parameter werden in der Datei \texttt{parameters.dat}
zusammengefasst (vgl.\ Tabelle~\ref{tab:evax-params}).
Die Option \texttt{-autoES} veranlasst \evax{}, vor der eigentlichen
Optimierung die Werte von $S_0^2$ und $\Delta E_0$ automatisch
zu bestimmen.
Anschließend durchläuft der Algorithmus eine Screening-Phase, in der
die grobe Strukturoptimierung stattfindet, gefolgt von einer
Feinoptimierung bis zum Erreichen des \texttt{Froze\_in}-Kriteriums.

\subsection{Bewertung der Fitgüte}

Die Güte des Fits wird durch den Residuenwert $R$ quantifiziert, der
als normierte Abweichung zwischen berechnetem und experimentellem
Signal im gewählten Vergleichsraum definiert ist.
Werte unter $R \approx 0{,}1$ gelten als guter Fit; bei Multi-Edge-Fits
ist zu beachten, dass der gewichtete Gesamtresidual die
Einzelresiduen der Kanten zusammenfasst und daher höher ausfallen kann.

Neben dem numerischen Residual ist die visuelle Inspektion der
$\chi(k)$- und FT-Spektren unverzichtbar:
Systematische Abweichungen in bestimmten $k$- oder $R$-Bereichen
können auf Modelldefizite hinweisen, die im Residuenwert allein
nicht sichtbar werden.

\subsection{Interpretation der Strukturparameter}

Aus dem finalen Strukturmodell lassen sich berechnen:
\begin{itemize}[nosep]
  \item Partielle radiale Verteilungsfunktionen $g_{ij}(R)$,
  \item Koordinationszahlen $N_{ij}$ durch Integration des ersten
        Peaks von $g_{ij}(R)$,
  \item Mittlere Bindungslängen $\langle R_{ij}\rangle$ und deren
        Standardabweichungen,
  \item Warren-Cowley-Parameter $\alpha_{ij}$
        (vgl.\ Abschnitt~\ref{sec:sro}).
\end{itemize}

\begin{figure}[htbp]
  \centering
  \begin{tikzpicture}[
    node distance=1.2cm and 2cm,
    block/.style={
      draw, rounded corners, text width=5cm, align=center,
      minimum height=1cm, font=\small
    },
    arrow/.style={-{Stealth[length=6pt]}, thick}
  ]
    \node[block, fill=blue!10] (raw)
      {Rohdaten $\mu(E)$\\[0.2em]\scriptsize Athena / larch};
    \node[block, fill=blue!10, below=of raw] (chi)
      {$\chi(k)$-Extraktion\\[0.2em]\scriptsize Normierung, Hintergrund};
    \node[block, fill=green!10, below=of chi] (model)
      {Startmodell (Superzelle)\\[0.2em]\scriptsize POSCAR, $4\times4\times4$};
    \node[block, fill=orange!10, below=of model] (evax)
      {\evax{}-Optimierung\\[0.2em]\scriptsize EA + RMC, FEFF};
    \node[block, fill=red!10, below=of evax] (result)
      {Ergebnisse\\[0.2em]\scriptsize RDF, $N_{ij}$, $\alpha_{ij}$};

    \draw[arrow] (raw) -- (chi);
    \draw[arrow] (chi) -- (model);
    \draw[arrow] (model) -- (evax);
    \draw[arrow] (evax) -- (result);

    % Feedback loop
    \draw[arrow, dashed, gray] (result.east) -- ++(1.5,0)
      |- node[right, pos=0.25, font=\scriptsize\itshape, text=gray]
        {Parameter anpassen} (evax.east);
  \end{tikzpicture}
  \caption{Schematischer Workflow einer \evax{}-Analyse, von der
    Rohdatenverarbeitung bis zur strukturellen Interpretation.}
  \label{fig:workflow}
\end{figure}

% Vollständige Pipeline-Skizze
% =============================================================================
% Analyse-Pipeline: saubere vertikale Darstellung
% =============================================================================
\begin{figure}[p]
  \centering
  \begin{tikzpicture}[
    scale=0.82, transform shape,
    % Styles
    step/.style={
      draw, rounded corners=4pt, fill=#1,
      text width=12cm, align=left, minimum height=1cm,
      font=\small, inner sep=8pt
    },
    step/.default={blue!6},
    filebox/.style={
      draw, fill=yellow!12, rounded corners=2pt,
      font=\scriptsize\ttfamily, inner sep=4pt
    },
    heading/.style={
      font=\small\bfseries, anchor=west
    },
    arr/.style={-{Stealth[length=5pt]}, thick, gray!70},
    loopback/.style={-{Stealth[length=5pt]}, thick, red!50!black, rounded corners=4pt},
  ]

    % ── 1. Messung ──────────────────────────────────────────────────────────
    \node[step=blue!8] (S1) at (0, 0) {
      \textbf{1 \enspace Synchrotron-Messung}\\[3pt]
      Absorptionsspektrum $\mu(E)$ im Bereich
      \num{7500}--\SI{10000}{eV}.
      Ein einziger Energiescan deckt die K-Kanten
      von Co (\SI{7709}{eV}), Ni (\SI{8333}{eV}) und
      Cu (\SI{8979}{eV}) gleichzeitig ab.
    };

    % ── 2. Athena ───────────────────────────────────────────────────────────
    \node[step=orange!8, below=0.45cm of S1] (S2) {
      \textbf{2 \enspace Vorverarbeitung in Athena}\\[3pt]
      Import der Rohdaten, visuelle Kontrolle,
      erste $E_0$-Bestimmung.
      Export als Athena-Projekt
      \,\tikz[baseline=-0.3ex]{\node[filebox]{MEA\_athena.prj};}
    };
    \draw[arr] (S1) -- (S2);

    % ── 3. Extraktion ──────────────────────────────────────────────────────
    \node[step=orange!8, below=0.45cm of S2] (S3) {
      \textbf{3 \enspace $\chi(k)$-Extraktion mit larch}\\[3pt]
      Neu-Extraktion nötig, da Multi-Edge-Normierung in
      Athena fehlerhafte Kantensprünge lieferte.
      Pro Kante wird ein separates Energiefenster gesetzt
      (Co: ${\leq}\SI{8233}{eV}$,
       Ni: \SIrange{8100}{8879}{eV},
       Cu: ${\geq}\SI{8700}{eV}$).\\[3pt]
      Schritte: $E_0$ bestimmen $\to$
      Pre-Edge abziehen $\to$
      auf $\Delta\mu_0 = 1$ normieren $\to$
      Hintergrund entfernen (\texttt{autobk}, $R_\text{bkg}=\SI{1.0}{\AA}$)
      $\to$ Umrechnung $E \to k$.\\[4pt]
      Ausgabe:\enspace
      \tikz[baseline=-0.3ex]{\node[filebox]{MEA\_Co\_chi.dat};}\enspace
      \tikz[baseline=-0.3ex]{\node[filebox]{MEA\_Ni\_chi.dat};}\enspace
      \tikz[baseline=-0.3ex]{\node[filebox]{MEA\_Cu\_chi.dat};}\\
      {\scriptsize Jeweils zwei Spalten: $k$ (\AA$^{-1}$) und $\chi(k)$.}
    };
    \draw[arr] (S2) -- (S3);

    % ── 4. Startmodell ─────────────────────────────────────────────────────
    \node[step=green!8, below=0.45cm of S3] (S4) {
      \textbf{4 \enspace Startmodell (Superzelle)}\\[3pt]
      Kristallstruktur wählen (FCC, BCC oder HCP) und
      $4{\times}4{\times}4$-Superzelle mit 256 Atomen
      (85\,Co, 85\,Ni, 86\,Cu, zufällig verteilt) erzeugen.
      Format: POSCAR mit Gittervektoren und fraktionellen
      Koordinaten.\\[4pt]
      Ausgabe:\enspace
      \tikz[baseline=-0.3ex]{\node[filebox]{MEA\_CoNiCu.p1};}
    };
    \draw[arr] (S3) -- (S4);

    % ── 5. Konfiguration ──────────────────────────────────────────────────
    \node[step=green!8, below=0.45cm of S4] (S5) {
      \textbf{5 \enspace Konfiguration}
      \hfill\tikz[baseline=-0.3ex]{\node[filebox]{parameters.dat};}\\[3pt]
      Zentrale Steuerdatei mit allen Einstellungen:
      $k$-Bereich (\SIrange{3}{11.5}{\per\AA}),
      $k$-Gewichtung ($k^2$),
      Vergleichsraum (Wavelet),
      Streupfade ($N_\text{legs}=4$),
      FEFF-Clusterradius (\SI{6}{\AA}),
      Schrittweite, Konvergenzkriterien.\\[3pt]
      Die Option \texttt{-autoES} bestimmt $S_0^2$ und $\Delta E_0$
      automatisch vor dem Hauptlauf.
    };
    \draw[arr] (S4) -- (S5);

    % ── 6. EvAX-Kern ──────────────────────────────────────────────────────
    \node[step=red!6, below=0.45cm of S5,
          draw=red!40!black, thick] (S6) {
      \textbf{6 \enspace EvAX-Optimierung}
      \hfill{\small\texttt{./EvAX.exe parameters.dat -autoES}}\\[4pt]
      %
      \begin{minipage}[t]{5.5cm}
        \textbf{Iterationsschleife:}\\[2pt]
        {\scriptsize
        \textbf{a)} Zufälliges Atom verschieben\\
        \quad\enspace (max.\ \SI{0.005}{\AA}/Schritt)\\[1pt]
        \textbf{b)} FEFF berechnet Streuphasen\\
        \quad\enspace für die neue Umgebung\\[1pt]
        \textbf{c)} $\chi_\text{calc}(k)$ mitteln über alle\\
        \quad\enspace Absorber gleichen Typs\\[1pt]
        \textbf{d)} Vergleich mit $\chi_\text{exp}(k)$\\
        \quad\enspace im Wavelet-Raum\\[1pt]
        \textbf{e)} Metropolis-Akzeptanz:\\
        \quad\enspace besser $\to$ behalten,\\
        \quad\enspace schlechter $\to$ evtl.\ behalten\\[1pt]
        $\to$ zurück zu \textbf{a)}
        }
      \end{minipage}%
      \hfill
      \begin{minipage}[t]{5.5cm}
        \textbf{FEFF-Dateien im Arbeitsverz.:}\\[2pt]
        {\scriptsize
        \tikz[baseline=-0.3ex]{\node[filebox]{feff};} Binary (feff85L)\\[2pt]
        \tikz[baseline=-0.3ex]{\node[filebox]{run-FEFF.exe};} Shell-Wrapper\\[2pt]
        \tikz[baseline=-0.3ex]{\node[filebox]{feff\_0/ feff\_1/};} je Kante:\\
        \quad \texttt{feff.inp} -- Eingabe (Cluster)\\
        \quad \texttt{feff.bin} -- Streuphasen\\
        \quad \texttt{chi.dat} -- berechnetes $\chi(k)$\\[4pt]
        \textbf{Evolutionärer Algorithmus:}\\[2pt]
        16 parallele Strukturkopien (States).\\
        Alle $\sim$100 Iterationen: schlechte\\
        States durch gute ersetzen +\\
        leichte Mutationen.
        }
      \end{minipage}
    };
    \draw[arr] (S5) -- (S6);

    % ── 7. Ausgaben ───────────────────────────────────────────────────────
    \node[step=purple!6, below=0.45cm of S6] (S7) {
      \textbf{7 \enspace Ausgaben}
      \hfill{\scriptsize Verzeichnis \texttt{output\_files/}}\\[4pt]
      \tikz[baseline=-0.3ex]{\node[filebox]{output.dat};}
        {\scriptsize\enspace Residual vs.\ Iteration (Konvergenz)}\hfill
      \tikz[baseline=-0.3ex]{\node[filebox]{final.xyz};}
        {\scriptsize\enspace optimierte 256 Atompositionen}\\[3pt]
      \tikz[baseline=-0.3ex]{\node[filebox]{Co/EXAFS.dat};}
        {\scriptsize\enspace $\chi_\text{calc}(k)$ und $\chi_\text{exp}(k)$
         pro Kante}\hfill
      \tikz[baseline=-0.3ex]{\node[filebox]{Co/expFT.dat};}
        {\scriptsize\enspace Fouriertransformierte}\\[3pt]
      \tikz[baseline=-0.3ex]{\node[filebox]{Co/RDF\_Co\_Ni.dat};}
        {\scriptsize\enspace partielle Paarverteilungsfunktionen $g_{ij}(R)$}\hfill
      \tikz[baseline=-0.3ex]{\node[filebox]{restart*};}
        {\scriptsize\enspace Checkpoint zum Fortsetzen}
    };
    \draw[arr] (S6) -- (S7);

    % ── 8. Auswertung ────────────────────────────────────────────────────
    \node[step=green!10, below=0.45cm of S7,
          draw=green!40!black] (S8) {
      \textbf{8 \enspace Auswertung (Python)}\\[4pt]
      \begin{minipage}[t]{5.5cm}
        \tikz[baseline=-0.3ex]{\node[filebox]{plot\_results.py};}\\[2pt]
        {\scriptsize Parst \texttt{output\_files/},
         erzeugt Konvergenz-, EXAFS-,
         FT-, RDF- und SRO-Plots.}
      \end{minipage}%
      \hfill
      \begin{minipage}[t]{5.5cm}
        \tikz[baseline=-0.3ex]{\node[filebox]{parameter\_scan.py};}\\[2pt]
        {\scriptsize Automatisierte Variation aller
         Fit-Parameter (83 Kombinationen,
         6 parallele EvAX-Instanzen).}
      \end{minipage}\\[6pt]
      \textbf{Physikalische Ergebnisse:}\enspace
      {\scriptsize
      Koordinationszahlen $N_{ij}$,\enspace
      Bindungslängen $\langle R_{ij}\rangle$,\enspace
      Warren-Cowley-Parameter $\alpha_{ij}$,\enspace
      Kristallstruktur.}
    };
    \draw[arr] (S7) -- (S8);

  \end{tikzpicture}
  \caption{Vollständige Analyse-Pipeline von der Synchrotron-Messung
    bis zu den physikalischen Ergebnissen.
    Gelbe Kästen bezeichnen konkrete Dateien.
    Der rot umrandete Schritt~6 zeigt den iterativen Kern von \evax{},
    in dem FEFF-Rechnungen und RMC-Optimierung zusammenwirken.}
  \label{fig:full-pipeline}
\end{figure}


\section{Stärken, Grenzen und Fehlerquellen}

\subsection{Stärken}

\begin{itemize}[nosep]
  \item Simultane Auswertung beliebig vieler Kanten mit einem
        konsistenten Strukturmodell,
  \item Keine Parametrisierung der RDF nötig -- das Modell kann auch
        asymmetrische oder multimodale Verteilungen abbilden,
  \item Direkte Berechnung von SRO-Parametern aus der Struktur,
  \item Eingebauter Schutz vor lokalen Minima durch den EA.
\end{itemize}

\subsection{Grenzen}

\begin{itemize}[nosep]
  \item Die Rechenzeit skaliert mit der Zahl der Atome, der Kanten
        und der Berücksichtigung von MS-Pfaden; ein Lauf mit
        256~Atomen und drei Kanten benötigt typischerweise mehrere
        Stunden,
  \item Die Topologie des Gitters ist durch das Startmodell
        vorgegeben und wird nicht optimiert,
  \item Bei geringer Datenqualität (niedriges Signal-zu-Rausch-Verhältnis,
        kleiner nutzbarer $k$-Bereich) ist die Struktur
        unterbestimmt, und der Algorithmus kann verschiedene
        gleichwertige Minima finden.
\end{itemize}

\subsection{Typische Fehlerquellen}

\begin{itemize}[nosep]
  \item Fehlerhafte $\chi(k)$-Extraktion (falsche $E_0$-Zuordnung,
        unzureichende Normierung),
  \item Ungeeigneter $k$-Bereich (zu viel Rauschen bei hohem $k$,
        Hintergrundartefakte bei niedrigem $k$),
  \item Falsche Stöchiometrie oder Kristallstruktur im Startmodell,
  \item Unzureichende Konvergenz (\texttt{Froze\_in} zu klein gewählt).
\end{itemize}
